% Appendix B. Auxiliary results: every external theorem the proofs of Appendix C use, in the form in which it is used.
% Sources: theory/framework.md and theory/propositions.json (proofs, repaired proofs), theory/round2_verdicts.json
% (suggested fixes naming the constant-rank theorem, transversality, Frobenius, Segre, Eckart--Young, zero sets),
% site/sections/13-appendix.html (Parts C and D) and site/sections/04-erlangen.html (the LOFT example).
% Numbering B.1--B.39; each entry ends with the results that use it. Bibliography: bib/extra/B-auxiliary.bib.
\providecommand{\fbBusedin}[1]{\par\smallskip\noindent{\small\color{inkgray}{\sffamily\bfseries\color{tide}Used in.}\enspace #1\par}\medskip}

\section{Auxiliary results}\label{app:auxiliary}

This appendix quotes the results from outside the paper that the proofs in Appendix~\ref{app:proofs} use. Each is
stated in the form in which it is applied, with a precise source, and is followed by the places that use it. Several are
textbook facts that a specialist applies without comment. They are listed so that every step of Appendix~\ref{app:proofs}
can be traced to a statement with hypotheses. For a few elementary facts we give the one-line reason instead of a
reference. Background definitions (manifolds, germs, group actions, lenses) are in Appendix~\ref{app:background}, and
where an entry has a companion there we name it.

Three results of the main text are external results restated in the paper's conventions and proved there:
\thmref{Theorem}{II.10} (the min-cut bound, here \auxref{B.14}), \thmref{Theorem}{III.5} (Noether charges,
\auxref{B.29}) and \thmref{Proposition}{III.7}(a) (the Kempf--Ness slice, \auxref{B.26}). The entries below give their
source forms.

The entries are grouped by subject: smooth maps and limits (B.1--B.4), real algebraic and tame geometry (B.5--B.10),
matrix analysis (B.11--B.21), Lie groups and invariant theory (B.22--B.27), dynamics, optimisation and differentiation
(B.28--B.35), three elementary facts (B.36), and spectral projectors, Riemannian submanifolds and floating-point
scaling (B.37--B.39). Notation follows the exposés: $\Mr$ is the set of $m\times n$ matrices of
rank at most $r$, $\MM_r$ those of rank exactly $r$, $P_U$ the orthogonal projector onto a subspace $U$, and
$\lVert\cdot\rVert$ the Frobenius norm unless marked otherwise. Entries of Appendix~\ref{app:background} are cited as
``Background A.1'' to ``Background D.34''; the numbers of its Part~B are unrelated to the numbers B.1--B.39 used here.

% ================================================================================================
\subsubsection{Smooth maps, fibre products and limits}

\begin{auxresult}{B.1}{rank, the rank theorem and submersions}{\citep[Prop.~4.1, Thm~4.12, Prop.~4.28,
Thm~5.12]{xlee2013introduction}}
Let $F:Q\to\Theta$ be a smooth map between manifolds of dimensions $N$ and $D$.
\begin{enumerate}[label=(\alph*)]
  \item For every $k$, the set $\{q:\rk dF_q\ge k\}$ is open. In particular the set where $\rk dF_q$ takes its maximal
    value is open and nonempty, and if $\rk dF_p=D$ then $F$ is a submersion on a neighbourhood of $p$.
  \item (Rank theorem.) If $\rk dF_q=k$ for all $q$ in a neighbourhood of $p$, there are charts centred at $p$ and at
    $F(p)$ in which $F$ reads $(x_1,\dots,x_N)\mapsto(x_1,\dots,x_k,0,\dots,0)$. Hence $p$ has a neighbourhood $V$ such
    that $F(V)$ is an embedded submanifold of dimension $k$ and $F^{-1}(F(p))\cap V$ is an embedded submanifold of
    dimension $N-k$.
  \item A submersion is an open map. In particular, if $\rk dF=D$ on a neighbourhood $U$ of $p$, then $F(U)$ contains
    an open neighbourhood of $F(p)$.
\end{enumerate}
\end{auxresult}
\noindent\emph{Reason for (a).} In coordinates, $\rk dF_q\ge k$ says that some $k\times k$ minor of the Jacobian is
nonzero at $q$, an open condition; Lee's Prop.~4.1 is the case $k=D$. Companion: \bgref{C.5}.
\fbBusedin{\thmref{Observation}{I.4}(b), where the lifts of both axes force $\rk d\iota_{p^*}=2$, the rank stays $2$
near $p^*$, and the image of a neighbourhood of $p^*$ contains an open subset of $\R^2$; \thmref{Definition}{I.5}
(with \auxref{B.9}); \thmref{Proposition}{II.12}(a) (fibre dimension $\lvert M\rvert-d(M)$ near points of maximal
rank); \thmref{Proposition}{III.4}(d) (rank $r$ is an open condition on the slice). The proofs of \thmref{Proposition}{I.8}, \thmref{Theorem}{II.7} and \thmref{Proposition}{IV.3} in Appendix~\ref{app:proofs} also cite it.}

\begin{auxresult}{B.2}{transversality and fibre products}{\citep[Cor.~5.30 and Thm~6.30]{xlee2013introduction};
\citep[\S1.5]{xguillemin1974differential}}
Let $\rho_M:Q_M\to\Theta$ and $\rho_N:Q_N\to\Theta$ be smooth, with $\dim Q_M=\lvert M\rvert$ and
$\dim Q_N=\lvert N\rvert$, and let $P=\{(q,q'):\rho_M(q)=\rho_N(q')\}\subseteq Q_M\times Q_N$.
\begin{enumerate}[label=(\alph*)]
  \item If $d\rho_M(T_qQ_M)+d\rho_N(T_{q'}Q_N)=T_\theta\Theta$ at every $(q,q')\in P$, then $P$ is an embedded
    submanifold of $Q_M\times Q_N$ of dimension $\lvert M\rvert+\lvert N\rvert-\dim\Theta$.
  \item Transversality at a single point of $P$ forces $\lvert M\rvert+\lvert N\rvert\ge\dim\Theta$.
  \item If $S\subseteq X$ is an embedded submanifold, every smooth map $T\to X$ with image in $S$ is smooth as a map
    into $S$. Consequently, whenever $P$ is an embedded submanifold of $Q_M\times Q_N$ (under (a), or when $\rho_M$ and
    $\rho_N$ are affine, so that $P$ is an affine subspace), $P$ with its two projections is the pullback of $\rho_M$ and
    $\rho_N$ in $\Diff$.
\end{enumerate}
\end{auxresult}
\noindent\emph{Reason.} $P=(\rho_M\times\rho_N)^{-1}(\Delta_\Theta)$, and $\rho_M\times\rho_N$ is transverse to the
diagonal $\Delta_\Theta$, of codimension $\dim\Theta$, exactly when the condition in (a) holds; Lee's Thm~6.30(a) gives
(a). Part (b) is a dimension count, since the left side of the condition has dimension at most
$\lvert M\rvert+\lvert N\rvert$. Companion: \bgref{C.7}.
\fbBusedin{\thmref{Observation}{I.4}(b), \thmref{Proposition}{I.8}(a), \thmref{Theorem}{II.7}(d). The proof of \thmref{Proposition}{III.4} in Appendix~\ref{app:proofs} also cites it.}

\begin{auxresult}{B.3}{limits in slices and coslices; right adjoints}{\citep[\S III.4, \S V.1 Exercise~1, \S V.5
Thm~1]{xmaclane1998categories}}
Let $\mathcal C$ be a category.
\begin{enumerate}[label=(\alph*)]
  \item In the slice $\mathcal C/X$, a product of $a:A\to X$ and $b:B\to X$ is a pullback $A\times_XB$ of $a$ and $b$
    in $\mathcal C$, with its map to $X$; one exists exactly when the other does.
  \item For an object $X$, the projection from the coslice $X/\mathcal C$ (the comma category $(X\downarrow\mathcal C)$)
    to $\mathcal C$ creates limits: every limit cone in $\mathcal C$ of the underlying diagram lifts to exactly one cone
    in $X/\mathcal C$ over the given diagram, and that cone is a limit.
  \item A functor that has a left adjoint preserves all limits that exist in its domain. The forgetful functors
    $\Set_*\to\Set$ and $\Diff_*\to\Diff$ have the left adjoint $X\mapsto X\sqcup\mathrm{pt}$.
\end{enumerate}
\end{auxresult}
\noindent\emph{Form used.} A product in the pointed slice $\Meth(\Theta,\theta_0)=\Diff_*/(\Theta,\theta_0)$ is a
pullback over $(\Theta,\theta_0)$ in $\Diff_*$ by (a). By (b) a pullback in $\Diff$ of the underlying maps lifts to
$\Diff_*$ with base point $(q_0,q_0')$, and by (c) any pullback in $\Diff_*$ is a pullback in $\Diff$, so it may be
tested with unpointed curves. Companions: \bgref{A.11}, \bgref{A.12}, \bgref{A.13}.
\fbBusedin{\thmref{Observation}{I.4}(b) (products in $\Meth$ as pullbacks, tested with unpointed curves). The proof of \thmref{Proposition}{I.8} in Appendix~\ref{app:proofs} also cites it.}

\begin{auxresult}{B.4}{smooth descent along submersions; quotient manifold theorem}{\citep[Thms~4.29, 4.30,
21.10]{xlee2013introduction}}
\begin{enumerate}[label=(\alph*)]
  \item Let $\pi:X\to Y$ be a surjective smooth submersion. A map $\bar f:Y\to Z$ is smooth iff $\bar f\circ\pi$ is
    smooth, and every smooth $f:X\to Z$ that is constant on the fibres of $\pi$ factors uniquely as $f=\bar f\circ\pi$
    with $\bar f$ smooth.
  \item If a Lie group $G$ acts smoothly, freely and properly on a manifold $X$, then $X/G$ has a unique smooth structure
    for which $\pi:X\to X/G$ is a submersion, and $\dim X/G=\dim X-\dim G$.
\end{enumerate}
\end{auxresult}
\noindent\emph{Instance.} On pairs $(B,A)$ of full rank, $(B,A)\mapsto BA$ is a surjective submersion onto $\MM_r$ whose
fibres are the $\GL_r$-orbits $\{(Bg^{-1},gA)\}$ (\auxref{B.5}(b)), so a function of the factors that is constant on
these orbits is a smooth function of $W$ on the full-rank locus. Companion: \bgref{D.22}.
\fbBusedin{\thmref{Proposition}{III.4}(d) (cometrics that are functions of $W$), \thmref{Proposition}{V.4}(a)
(descent of combination rules on the full-rank locus).}

% ================================================================================================
\subsubsection{Real algebraic and tame geometry}

\begin{auxresult}{B.5}{determinantal varieties}{\citep[Lectures~9, 12 and 14]{xharris1992algebraic};
\citep[Thm~3.2]{xschneider2015convergence}}
Let $1\le r<\min(m,n)$.
\begin{enumerate}[label=(\alph*)]
  \item $\Mr\subseteq\R^{m\times n}$ is the common zero set of the $(r+1)\times(r+1)$ minors. It is closed and closed
    under all real scalings, so it is a closed cone with vertex $0$, of dimension $r(m+n-r)$.
  \item $\MM_r$ is open and dense in $\Mr$ and is an embedded submanifold of dimension $r(m+n-r)$. If $X\in\MM_r$ has
    column space $U$ and row space $V$, and $X=BA$ with $B\in\R^{m\times r}$, $A\in\R^{r\times n}$, then
    \begin{equation*}
      T_X\MM_r=\{\dot BA+B\dot A\}=\{D:P_{U^\perp}DP_{V^\perp}=0\},
    \end{equation*}
    and $(B,A)\mapsto BA$ is a submersion from the pairs of full rank onto $\MM_r$ whose fibres are the orbits
    $\{(Bg^{-1},gA):g\in\GL_r\}$.
  \item At a point $X$ of rank $s<r$, with column space $U$ and row space $V$, the tangent cone of $\Mr$ is
    $\{D:\rk(P_{U^\perp}DP_{V^\perp})\le r-s\}$, which is not a linear space. Hence $\Mr$ is not a $C^1$ submanifold near
    $X$, and $\mathrm{Sing}\,\Mr=\MM_{\le r-1}$, algebraically and in the $C^1$ sense.
\end{enumerate}
\end{auxresult}
\noindent\emph{Reason for the last claim of (c).} $P_{U^\perp}DP_{V^\perp}$ ranges over all $(m-s)\times(n-s)$
matrices, and since $r-s<\min(m-s,n-s)$ two of rank $r-s$ can sum to a matrix of larger rank. The tangent cone of a
$C^1$ submanifold is its tangent space, a linear space. Companion: \bgref{C.12}.
\fbBusedin{\thmref{Theorem}{II.1}(c), \thmref{Theorem}{II.2}(a), (b), (c), (e) (the image determines its vertex),
\thmref{Proposition}{II.3}, \thmref{Proposition}{II.8} (tangent space at $(I,X)$), \thmref{Proposition}{II.12},
\thmref{Proposition}{II.16} (the translation stabiliser of $\Mr$ is $\{0\}$), \thmref{Proposition}{III.4}(b), (d),
\thmref{Proposition}{V.4}(a). The proofs of \thmref{Theorem}{II.7} and \thmref{Proposition}{II.9} in Appendix~\ref{app:proofs} also cite it.}

\begin{auxresult}{B.6}{cones that are $C^1$ at their vertex}{elementary}
Let $C\subseteq\R^N$ be a closed set with $0\in C$ and $tC\subseteq C$ for every $t>0$. If $C$ is a $C^1$ embedded
submanifold near $0$, then $C$ equals the linear subspace $T_0C$. Consequently, the vertex of a closed cone that is not a
linear subspace is a singular point.
\end{auxresult}
\noindent\emph{Reason.} The tangent cone of $C$ at $0$, the set of limits of $t_kx_k$ with $x_k\in C$, $x_k\to0$ and
$t_k>0$, equals $C$ because $C$ is a closed cone, and the tangent cone of a $C^1$ submanifold at a point is its tangent
space. Companions: \bgref{C.9}, \bgref{C.10}.
\fbBusedin{\thmref{Definition}{I.5}, \thmref{Theorem}{II.2}(a), (c), \thmref{Theorem}{II.7}(c),
\thmref{Theorem}{II.11}(e) (the closed cone $\mathcal K_{\le k}$).}

\begin{auxresult}{B.7}{real normal forms; the skew-determinantal cone}{\citep[Ch.~2]{xhorn2013matrix}}
\begin{enumerate}[label=(\alph*)]
  \item Every $\Omega\in\so(n)$ is orthogonally similar to a block-diagonal matrix with $2\times2$ blocks
    $\mu_k\left(\begin{smallmatrix}0&-1\\1&0\end{smallmatrix}\right)$, $\mu_1\ge\mu_2\ge\dots>0$, and zeros. Hence
    $\rk\Omega$ is even, the nonzero eigenvalues are $\pm i\mu_k$, the singular values are $\mu_1,\mu_1,\mu_2,\mu_2,\dots$,
    and $\Omega^4$ is symmetric positive semidefinite with eigenvalues $\mu_k^4$. A skew matrix of odd size has $0$ as an
    eigenvalue.
  \item Every $R\in\SO(n)$ is orthogonally similar to a block-diagonal matrix with $2\times2$ rotation blocks of angles
    $\varphi_k\in(0,\pi]$ and ones. Then $\rk(R-I)$ is twice the number of blocks, $\mathrm{im}(R-I)$ is the sum of the
    rotation planes, and $R=e^X$ for a skew $X$ whose eigenvalues are $\pm i\varphi_k$ and $0$.
  \item For even $r$ with $2\le r\le n-2$, the set $\Sigma_r=\{\Omega\in\so(n):\rk\Omega\le r\}$ is a closed cone of
    dimension $r(n-r)+r(r-1)/2=rn-r(r+1)/2$. Its points of rank exactly $r$ form an open dense submanifold of that
    dimension, and $\Sigma_r$ is not a linear space.
\end{enumerate}
\end{auxresult}
\noindent\emph{Reason for (c).} A skew matrix of rank $r$ is a range $V\in\Gr(r,n)$ together with a nondegenerate skew
form on $V$, an open subset of $\so(r)$; and skew matrices of ranks $r$ and $2$ can sum to rank $r+2\le n$. Companion:
\bgref{C.13}.
\fbBusedin{\thmref{Corollary}{II.6} (the Cayley--Neumann factor), \thmref{Theorem}{II.7}(a), (c), (d) (the dimension count
and the density argument), \thmref{Theorem}{V.3}, corollary~4, the LOFT example of Exposé~II.B.}

\begin{auxresult}{B.8}{Segre cones}{\citep[Lecture~2]{xharris1992algebraic}}
The Segre cone $\MM_{\le1}(p\times q)=\{ab\T:a\in\R^p,b\in\R^q\}$ is a closed cone of dimension $p+q-1$, it is not a
linear space when $p,q\ge2$, and it spans $\R^{p\times q}$. Hence $k$ generic rank-one matrices span a space of
dimension $\min(k,pq)$: the parameters $(a_1,b_1,\dots,a_k,b_k)$ for which $a_1b_1\T,\dots,a_kb_k\T$ span less form a
proper algebraic subset of $(\R^p\times\R^q)^k$.
\end{auxresult}
\noindent\emph{Reason for the last claim.} Spanning dimension at least $d$ is the nonvanishing of some $d\times d$
minor of the coefficient matrix, a polynomial in the parameters, and distinct matrix units $E_{ij}$ show that it is not
identically zero for $d=\min(k,pq)$. Companion: \bgref{C.14}.
\fbBusedin{\thmref{Proposition}{II.8} (the $m$ points $(E\T e_i)\otimes(F\T w_{0,i})$ give the generic value of
$d(\mathrm{DoRA}_r)$), \thmref{Proposition}{II.12} (the KronA image), the VeRA card of Exposé~VII.}

\begin{auxresult}{B.9}{semialgebraic and definable sets; dimension}{\citep[Ch.~2 and 9]{xbochnak1998real};
\citep[Ch.~4, 6 and 7]{xvandendries1998tame}; \citep{xwilkie1996model}}
\begin{enumerate}[label=(\alph*)]
  \item (Tarski--Seidenberg.) The image of a semialgebraic set under a semialgebraic map is semialgebraic.
    (Wilkie.) The expansion $\R_{\exp}$ of the real field by the exponential function is o-minimal, so sets and maps
    defined by first-order formulas from polynomials and $\exp$ (softmax included) are definable, and images of
    definable sets under definable maps are definable.
  \item Every definable set $X$ is a finite disjoint union of $C^1$ embedded submanifolds (cells), and $\dim X$, the
    largest of their dimensions, does not depend on the decomposition. If $X\subseteq Y$ then $\dim X\le\dim Y$, and
    $\dim f(X)\le\dim X$ for every definable map $f$.
  \item (Definable choice and generic smoothness.) A definable surjection has a definable section, and a definable map
    is $C^1$ outside a definable subset of its domain of smaller dimension.
  \item Consequently, for a definable $C^1$ map $\rho$ on an open definable set $Q\subseteq\R^k$,
    \begin{equation*}
      \dim\rho(Q)=\max_{q\in Q}\rk d\rho_q .
    \end{equation*}
\end{enumerate}
\end{auxresult}
\noindent\emph{Reason for (d).} At a point of maximal rank, the rank theorem (\auxref{B.1}(b)) puts a sheet of that
dimension inside $\rho(Q)$. Conversely, by (c) a definable section $s$ of $\rho$ is $C^1$ on an open subset $S$ of a
top-dimensional cell of $\rho(Q)$, and $\rho\circ s=\id_S$ gives $\rk d\rho_{s(y)}\ge\dim S=\dim\rho(Q)$. Companions:
\bgref{C.17}, \bgref{C.18}.
\fbBusedin{\thmref{Definition}{I.5} ($d(M)=\dim\Im M=\max_q\rk d\rho_q$), the obstruction $d(M)>d(N)$ of
Exposé~VIII.C, \thmref{Proposition}{II.12}(a). The proofs of \thmref{Theorem}{II.7}, \thmref{Proposition}{II.8}, \thmref{Proposition}{II.9}, \thmref{Proposition}{II.13} and \thmref{Theorem}{V.3} in Appendix~\ref{app:proofs} also cite it.}

\begin{auxresult}{B.10}{genericity; zero sets; the identity theorem}{\citep{xmityagin2020zero};
\citep[Ch.~1 and 2]{xkrantz2002primer}}
\begin{enumerate}[label=(\alph*)]
  \item Let $U\subseteq\R^N$ be open and connected and $f:U\to\R$ real-analytic, for instance a polynomial. If $f$
    vanishes on a nonempty open subset of $U$, then $f\equiv0$. If $f\not\equiv0$, its zero set is closed, nowhere dense
    and of Lebesgue measure zero.
  \item Hence a property that fails only on the common zero set of finitely many real-analytic functions on a connected
    domain, not all identically zero, holds on an open dense set of full measure, and with probability one under every
    law with a density. To prove such genericity it suffices to write failure as the vanishing of these functions and to
    exhibit one point where one of them does not vanish. A countable union of such exceptional sets still has measure
    zero.
  \item A function equal to a convergent power series $\sum_ka_kt^k$ near $0$ determines its coefficients. In
    particular, if $\sigma(ct)=c\,\sigma(t)$ near $0$ for some $c\ne0$, then $ca_k=c^ka_k$ for every $k$.
\end{enumerate}
\end{auxresult}
\noindent Companions: \bgref{C.21}, \bgref{C.22}.
\fbBusedin{\thmref{Theorem}{II.2}(d) (a generic rank-one update), \thmref{Proposition}{II.8}, \thmref{Theorem}{II.11}(c)
(generic Schmidt ranks), \thmref{Proposition}{II.12} (the generic rank of GraLoRA), \thmref{Proposition}{V.2}(a) ($e$
has full rank for almost every $W$, cell by cell of the quantiser), \thmref{Lemma}{VI.2}(c), \thmref{Theorem}{VI.3}
(the GELU obstruction), \thmref{Theorem}{VI.5}(f) (prefixes with affine output form a nowhere-dense analytic subset). The proofs of \thmref{Proposition}{I.8}, \thmref{Theorem}{II.1}, \thmref{Corollary}{II.6}, \thmref{Theorem}{II.7}, \thmref{Proposition}{II.13}, \thmref{Proposition}{IV.3} and \thmref{Theorem}{V.3} in Appendix~\ref{app:proofs} also cite it.}

% ================================================================================================
\subsubsection{Matrix analysis}

\begin{auxresult}{B.11}{singular value decomposition; Eckart--Young--Mirsky}{\citep{xeckart1936approximation,
xmirsky1960symmetric}; \citep[Ch.~2 and 7]{xhorn2013matrix}}
\begin{enumerate}[label=(\alph*)]
  \item Every $X\in\R^{m\times n}$ of rank $k$ has a singular value decomposition $X=\sum_{i\le k}\sigma_iu_iv_i\T$ with
    $\sigma_1\ge\dots\ge\sigma_k>0$ and orthonormal families $(u_i)$, $(v_i)$; two matrices lie in one orbit of
    $(P,R)\mapsto PXR$, $P\in\Ort(m)$, $R\in\Ort(n)$, iff they have the same singular values. Applied to the
    $m_1\times m_2$ reshape of $u\in\R^{m_1}\otimes\R^{m_2}$ ($x\otimes y\mapsto xy\T$), the decomposition is the Schmidt
    decomposition $u=\sum_i\sigma_ix_i\otimes y_i$, whose number of terms, the Schmidt rank, is the rank of the reshape.
    Every matrix of rank at most $r$ factors as $QV$ with $Q\in\R^{m\times r}$ having orthonormal columns.
  \item (Eckart--Young--Mirsky.) For $0\le s<k$ the truncation $[X]_s=\sum_{i\le s}\sigma_iu_iv_i\T$ minimises
    $\lVert X-Y\rVert$ over $\rk Y\le s$, with minimum $\big(\sum_{i>s}\sigma_i^2\big)^{1/2}$, and the minimiser is unique
    iff $\sigma_s>\sigma_{s+1}$. The matrices $[X]_s$ and $X-[X]_s$ are Frobenius-orthogonal, so for every real $c$
    \begin{equation*}
      \lVert X-c[X]_s\rVert^2=\lVert X-[X]_s\rVert^2+(1-c)^2\lVert[X]_s\rVert^2 .
    \end{equation*}
  \item (Reduced-rank regression.) For $L\in\GL_n$, $\arg\min_{\rk Y\le s}\lVert(X-Y)L\rVert=[XL]_sL^{-1}$, unique iff
    $\sigma_s(XL)>\sigma_{s+1}(XL)$; and $[XL]_s=\Pi_sXL$, where $\Pi_s$ is the orthogonal projector onto the top-$s$
    eigenspace of $XLL\T X\T$.
\end{enumerate}
\end{auxresult}
\noindent\emph{Reason for (c).} Substitute $Z=YL$, which has the same rank as $Y$, and apply (b) to $XL$; the left
singular vectors of $XL$ are the eigenvectors of $XLL\T X\T$.
\fbBusedin{\thmref{Proposition}{I.8}(a) (the SVD frames of LoRA-XS), \thmref{Theorem}{II.1}(c) (splitting an SVD),
\thmref{Theorem}{II.5}(3), \thmref{Theorem}{II.11}(c) (Schmidt decompositions), \thmref{Proposition}{II.16}(a) and the
lemma in (b), \thmref{Proposition}{III.7}(a) (the balanced representative), \thmref{Theorem}{III.9} (compact SVD),
\thmref{Proposition}{V.2}(b), (c) (LoftQ's rank-$r$ step and its defect), \thmref{Proposition}{VI.4}(c) (the factorisation
$R\T V$), Prediction~X.5 (the distance to $k$-term Kronecker sums). The proofs of \thmref{Proposition}{II.8} and \thmref{Theorem}{V.3} in Appendix~\ref{app:proofs} also cite it.}

\begin{auxresult}{B.12}{rank inequalities; the Hadamard rank bound}{\citep[Ch.~4 and 5]{xhorn1991topics}}
For matrices of compatible sizes:
\begin{enumerate}[label=(\alph*)]
  \item $\rk(X+Y)\le\rk X+\rk Y$ and $\rk(XY)\le\min(\rk X,\rk Y)$;
  \item $\rk(C\otimes D)=\rk C\cdot\rk D$;
  \item (Hadamard rank bound) $\rk(X\odot Y)\le\rk X\cdot\rk Y$. Explicitly, if $X=B_1A_1$ and $Y=B_2A_2$ with inner
    dimensions $r_1$ and $r_2$, then
    \begin{equation*}
      X\odot Y=(B_1\bullet B_2)\,(A_1\T\bullet A_2\T)\T ,
    \end{equation*}
    where row $i$ of the face-splitting product $B_1\bullet B_2\in\R^{m\times r_1r_2}$ is the Kronecker product of the
    rows $i$ of $B_1$ and $B_2$. In particular $W_0\odot(BA)=\sum_{k\le r}\diag(b_k)\,W_0\,\diag(a_k)$, summed over the
    columns $b_k$ of $B$ and the rows $a_k$ of $A$, has rank at most $r\rk W_0$.
\end{enumerate}
\end{auxresult}
\noindent\emph{Reason for (c).} $X\odot Y$ is the submatrix of $X\otimes Y$ on the rows $(i,i)$ and the columns
$(j,j)$, so (b) applies. Companion: \bgref{B.13}.
\fbBusedin{\thmref{Theorem}{II.1}(c), \thmref{Proposition}{II.9}(a) (HiRA), \thmref{Theorem}{II.10} (subadditivity),
\thmref{Theorem}{II.11}(b), \thmref{Proposition}{II.12} (LoHa, GraLoRA), \thmref{Proposition}{II.13} (two HiRA cycles),
\thmref{Theorem}{V.3}, corollaries~1 and~5, the face-splitting arrow of rule A3 (Exposé~VIII.C).}

\begin{auxresult}{B.13}{the Van Loan--Pitsianis rearrangement; operator-Schmidt decomposition}{\citep{loan1993approximation};
\citep{xvanloan2000ubiquitous}}
Let $m=m_1m_2$ and $n=n_1n_2$, cut $\Delta\in\R^{m\times n}$ into an $m_1\times n_1$ grid of $m_2\times n_2$ blocks
$\Delta_{ij}$, and let $\mathcal R(\Delta)\in\R^{m_1n_1\times m_2n_2}$ have the rows $\vect(\Delta_{ij})\T$, ordered with
$i$ running fastest.
\begin{enumerate}[label=(\alph*)]
  \item $\mathcal R$ is a linear bijection that permutes entries, hence a Frobenius isometry, and
    $\mathcal R(C\otimes D)=\vect C\,\vect D\T$ for $C\in\R^{m_1\times n_1}$ and $D\in\R^{m_2\times n_2}$.
  \item Hence $\Delta=\sum_{i\le k}C_i\otimes D_i$ for some $C_i,D_i$ iff $\rk\mathcal R(\Delta)\le k$. The Kronecker
    rank $\rk_\otimes\Delta$ equals $\rk\mathcal R(\Delta)\le\min(m_1n_1,m_2n_2)$, and the nearest sum of $k$ Kronecker
    products in Frobenius norm is $\mathcal R^{-1}\big([\mathcal R(\Delta)]_k\big)$ (\auxref{B.11}(b)).
  \item The SVD of $\mathcal R(\Delta)$ gives the operator-Schmidt decomposition $\Delta=\sum_{k\le K}s_kC_k\otimes D_k$
    with $s_k>0$, Frobenius-orthonormal families $(C_k)$, $(D_k)$ and $K=\rk_\otimes\Delta$. Conversely, any expansion of
    this form with positive coefficients and Frobenius-orthonormal families is one, so its number of terms is
    $\rk_\otimes\Delta$.
\end{enumerate}
\end{auxresult}
\noindent Companions: \bgref{B.11}, \bgref{B.12}.
\fbBusedin{\thmref{Theorem}{II.11}(a)--(e), \thmref{Proposition}{II.13} (the Kronecker family as
$\mathcal R^{-1}_*\mathrm{LoRA}_k$), Prediction~X.5. The proof of \thmref{Proposition}{II.12} in Appendix~\ref{app:proofs} also cites it.}

\begin{auxresult}{B.14}{the min-cut bound for tensor networks}{\citep{xcui2016quantum}}
Let $\Gamma$ be a tensor network evaluated in $(\FinVect,\otimes)$, every wire attached to a box, and let
$\mathcal I\sqcup\mathcal O$ be a partition of its boxes. Contracting each side separately factors the value
$\llbracket\Gamma\rrbracket$ through the tensor product of the wires that cross the partition and of the open legs
attached on the wrong side, so $\rk\llbracket\Gamma\rrbracket$ is at most the product $w$ of their dimensions, and at most $\min w$ over all
partitions. For fixed bond dimensions the largest rank over all choices of tensors can be strictly smaller than the
minimal cut weight; equality holds, for instance, when every bond dimension is a power of one integer.
\end{auxresult}
\noindent\emph{Form used.} \thmref{Theorem}{II.10} states the bound with the paper's conventions (bare strands, copy
spiders counted once, formal sums) and proves it; only the upper bound is used. Companion: \bgref{B.8}.
\fbBusedin{\thmref{Theorem}{II.10} and its consequences, \thmref{Proposition}{II.12} (maximum ranks), the
``high-rank PEFT'' remark of Exposé~II.C.}

\begin{auxresult}{B.15}{manifolds of fixed Tucker and tensor-train rank}{\citep{xkoch2010dynamical}
(Tucker); \citep{xholtz2012manifolds} (TT); \citep{xuschmajew2013geometry} (hierarchical Tucker)}
\begin{enumerate}[label=(\alph*)]
  \item (Tucker.) For an admissible multilinear rank $(r_1,\dots,r_d)$ of tensors in $\R^{n_1\times\dots\times n_d}$,
    the tensors of exactly that rank form an embedded submanifold of dimension $\prod_kr_k+\sum_kr_k(n_k-r_k)$. In the
    parametrisation by a core in $\R^{r_1\times\dots\times r_d}$ and factor matrices of full column rank, with
    $\prod_kr_k+\sum_kn_kr_k$ parameters, the fibres are the orbits of $\prod_k\GL_{r_k}$, so the dimension is the
    parameter count minus $\sum_kr_k^2$.
  \item (Tensor train.) For an admissible TT rank $(r_1,\dots,r_{d-1})$, with $r_0=r_d=1$, the tensors of exactly that
    rank form an embedded submanifold of dimension $\sum_{k=1}^dr_{k-1}n_kr_k-\sum_{k=1}^{d-1}r_k^2$, the number of core
    entries minus the dimension of the bond gauge $\prod_{k=1}^{d-1}\GL_{r_k}$.
\end{enumerate}
\end{auxresult}
\noindent Here admissible means realised by some tensor; at minimal ranks the bond gauge acts freely. Companion:
\bgref{B.9}.
\fbBusedin{\thmref{Proposition}{II.12}(b), row ``TT / Tucker''.}

\begin{auxresult}{B.16}{circulant matrices and the discrete Fourier transform}{\citep[Thm~3.1]{xgray2006toeplitz}}
\begin{enumerate}[label=(\alph*)]
  \item A circulant matrix with first column $w\in\R^b$ satisfies $\mathrm{circ}(w)=F^{-1}\diag(Fw)\,F$, with $F$ the
    discrete Fourier matrix. The powers $P^0,\dots,P^{b-1}$ of the cyclic shift are a basis of the circulants, with Gram
    matrix $bI$ in the Frobenius inner product.
  \item Over $\R$, since the eigenvalues of a real circulant at conjugate frequencies are conjugate, $\mathrm{circ}(w)$ is
    orthogonally similar to a block-diagonal matrix of $2\times2$ rotation--scaling blocks and $1\times1$ blocks at the
    self-conjugate frequencies.
  \item For a real coefficient $c$ and a frequency $(u,v)$, the real part of the inverse two-dimensional transform of
    $c\,e_{(u,v)}$ is the cosine pattern
    \begin{equation*}
      \mathrm{Re}\,\mathrm{ifft2}(c\,e_{(u,v)})_{j,l}=\tfrac{c}{mn}\cos\big(2\pi(uj/m+vl/n)\big).
    \end{equation*}
    It is the same for $(u,v)$ and $(-u,-v)$, and it is even under $(j,l)\mapsto(-j,-l)$, indices taken modulo $m$
    and $n$.
\end{enumerate}
\end{auxresult}
\fbBusedin{\thmref{Proposition}{II.12} (the FourierFT gauge from conjugate pairs), \thmref{Proposition}{II.13} (FourierFT
and C3A as masks in transported frames). The proof of \thmref{Theorem}{II.10} in Appendix~\ref{app:proofs} also cites it.}

\begin{auxresult}{B.17}{Ky Fan's maximum principle; interlacing}{\citep{xfan1949theorem}; \citep[Ch.~III]{xbhatia1997matrix};
\citep[Ch.~3]{xhorn1991topics}}
\begin{enumerate}[label=(\alph*)]
  \item (Ky Fan.) Let $S\in\R^{n\times n}$ be symmetric with eigenvalues $\lambda_1\ge\dots\ge\lambda_n$ and let
    $1\le k\le n$. Then
    \begin{equation*}
      \max\{\tr(PSP\T):P\in\R^{k\times n},\ PP\T=I_k\}=\lambda_1+\dots+\lambda_k ,
    \end{equation*}
    attained when the rows of $P$ span a top-$k$ eigenspace. For $S=X\T X$ this reads
    \begin{equation*}
      \max_P\lVert XP\T\rVert^2=\sigma_1(X)^2+\dots+\sigma_k(X)^2 .
    \end{equation*}
  \item (Interlacing for compressions.) If $P\in\R^{k\times n}$ and $Q\in\R^{l\times m}$ have orthonormal rows, then
    $\sigma_j(QXP\T)\le\sigma_j(X)$ for every $j$ and every $X\in\R^{m\times n}$; for symmetric $S$,
    $\lambda_j(PSP\T)\le\lambda_j(S)$.
  \item Consequently, if $S\in\so(n)$ has eigenvalues $\pm i\mu_1,\pm i\mu_2,\dots$ with $\mu_1\ge\mu_2\ge\dots$, then
    for every $P\in\R^{k\times n}$ with orthonormal rows
    \begin{equation*}
      \lVert PSP\T\rVert^2\le2\sum_{j\le\lfloor k/2\rfloor}\mu_j^2 ,
    \end{equation*}
    with equality when the row space of $P$ contains the invariant subspace of the top $\lfloor k/2\rfloor$ rotation
    pairs of $S$.
\end{enumerate}
\end{auxresult}
\noindent\emph{Reason for (c).} $PSP\T$ is skew of size $k$, so by \auxref{B.7}(a) its singular values are
$\nu_1,\nu_1,\nu_2,\nu_2,\dots$, with an extra $0$ when $k$ is odd. By (b), $\nu_j=\sigma_{2j-1}(PSP\T)\le
\sigma_{2j-1}(S)=\mu_j$, and $\lVert PSP\T\rVert^2=2\sum_{j\le\lfloor k/2\rfloor}\nu_j^2$. If the row space contains an
invariant subspace $W$ of the top pairs, then $W^\perp$ is invariant too, and the compression is $S|_W$ plus zeros.
Companion: \bgref{D.32}.
\fbBusedin{\thmref{Proposition}{II.3} (alignment: among $A_0$ with $A_0A_0\T=cI$ the rate $\lVert GA_0\T\rVert^2$ is
largest on the top-$r$ right singular space of $G$), the LOFT example of Exposé~II.B, part (d).}

\begin{auxresult}{B.18}{singular values of products; the Loewner order}{\citep[Ch.~3]{xhorn1991topics};
\citep[Ch.~III]{xbhatia1997matrix}; \citep[Ch.~7]{xhorn2013matrix}}
\begin{enumerate}[label=(\alph*)]
  \item (Horn.) $\sigma_{i+j-1}(XY)\le\sigma_i(X)\,\sigma_j(Y)$. In particular
    \begin{equation*}
      \sigma_i(XY)\le\sigma_i(X)\,\lVert Y\rVert_{\rm op},\qquad\sigma_i(XY)\le\lVert X\rVert_{\rm op}\,\sigma_i(Y),
    \end{equation*}
    and a product of contractions (matrices with $\lVert\cdot\rVert_{\rm op}\le1$) is a contraction.
  \item If $\lVert Y\rVert_{\rm op}\le1$, then $Y\T Y\preceq I$, and $WY\T YW\T\preceq WW\T$ for every $W$.
  \item (Weyl monotonicity.) If $S\preceq T$ are symmetric, then $\lambda_i(S)\le\lambda_i(T)$ for every $i$. With (b),
    $\sigma_i(WY\T)\le\sigma_i(W)$ whenever $\lVert Y\rVert_{\rm op}\le1$.
\end{enumerate}
\end{auxresult}
\fbBusedin{\thmref{Corollary}{II.6} (the Cayley--Neumann factor contracts for small $Q$), \thmref{Theorem}{V.3},
corollary~4 (each default-OFT merge decreases $K_{\rm out}$ in the Loewner order and lowers every singular value).}

\begin{auxresult}{B.19}{polar, Cartan and $KAK$ decompositions}{\citep[Ch.~2 and 7]{xhorn2013matrix}}
\begin{enumerate}[label=(\alph*)]
  \item (Polar.) Every $g\in\GL_r$ is uniquely $g=k\,e^S$ with $k\in\Ort(r)$ and $S\in\Sym_r$, and
    $(k,S)\mapsto k\,e^S$ is a diffeomorphism $\Ort(r)\times\Sym_r\to\GL_r$. Hence $g\mapsto g\T g=e^{2S}$ identifies the
    cosets $\Ort(r)g$ with the positive definite matrices, $\Ort(r)\backslash\GL_r\cong\Sym^+_r$.
  \item (Cartan.) $\gl_r=\so(r)\oplus\Sym_r$ through $\xi=\tfrac12(\xi-\xi\T)+\tfrac12(\xi+\xi\T)$, with
    $[\so,\so]\subseteq\so$, $[\so,\Sym]\subseteq\Sym$ and $[\Sym,\Sym]\subseteq\so$.
  \item ($KAK$.) Every $g\in\GL^+(b)$ is $k_1\diag(d)\,k_2$ with $k_1,k_2\in\SO(b)$ and $d\in\R^b_{>0}$.
\end{enumerate}
\end{auxresult}
\noindent\emph{Reason.} (a) combines the polar form $g=kP$, $P$ positive definite and unique, with the fact that
$\exp:\Sym_r\to\Sym_r^+$ is a diffeomorphism (spectral theorem). (c) is the SVD with signs absorbed: if both singular
frames have determinant $-1$, flip the first column of each. Companion: \bgref{D.27}.
\fbBusedin{\thmref{Theorem}{III.3} (gradient descent on LoRA as a family indexed by $\Ort(r)\backslash\GL_r$),
\thmref{Theorem}{III.5}(iii), \thmref{Proposition}{III.7}(a) (reduction to $h(t)=f(e^{t\xi}\cdot x)$ with $\xi$ symmetric),
\thmref{Theorem}{V.3}, corollary~4 (the monoid generated by Cayley--Neumann factors).}

\begin{auxresult}{B.20}{Sylvester equations}{\citep{xbhatia1997how}; \citep[Ch.~4]{xhorn1991topics}}
For $S\in\R^{p\times p}$ and $T\in\R^{q\times q}$, the map $X\mapsto SX+XT$ on $\R^{p\times q}$ has the eigenvalues
$\lambda_i(S)+\lambda_j(T)$. So $SX+XT=C$ has a unique solution for every $C$ iff $S$ and $-T$ have no common
eigenvalue. If $S$ and $T$ are symmetric positive definite, the solution is $X=\int_0^\infty e^{-tS}Ce^{-tT}\,dt$.
\end{auxresult}
\noindent Companion: \bgref{D.33}.
\fbBusedin{\thmref{Proposition}{III.4} (practice: LoRA-Pro's Sylvester-optimal free matrix exists and is unique on the
full-rank locus), \thmref{Theorem}{III.12} (published variants). The proof of \thmref{Proposition}{IV.3} in Appendix~\ref{app:proofs} also cites it.}

\begin{auxresult}{B.21}{the nuclear norm as a factorisation norm}{\citep[Lemma~1]{xsrebro2004maximum};
see also \citet{xrecht2010guaranteed}}
For $X\in\R^{m\times n}$ and every $r\ge\rk X$,
\begin{equation*}
  \lVert X\rVert_*=\min\Big\{\tfrac12\big(\lVert B\rVert^2+\lVert A\rVert^2\big):B\in\R^{m\times r},\ A\in\R^{r\times n},
  \ BA=X\Big\}=\min_{BA=X}\lVert B\rVert\,\lVert A\rVert ,
\end{equation*}
attained at $B=U\Sigma^{1/2}O$, $A=O\T\Sigma^{1/2}V\T$ for an SVD $X=U\Sigma V\T$ padded to inner size $r$ and any
$O\in\Ort(r)$.
\end{auxresult}
\fbBusedin{\thmref{Proposition}{III.7}(a) (the minimum value $2\lVert BA\rVert_*$ on a full-rank orbit) and its commentary
(weight decay on the factors as a nuclear-norm penalty).}

% ================================================================================================
\subsubsection{Lie groups and invariant theory}

\begin{auxresult}{B.22}{the Cartan--Dieudonné theorem in Scherk's sharp form}{\citep{xscherk1950decomposition};
\citep[Ch.~III]{xartin1957geometric}}
For $u\ne0$ let $H_u=I-2uu\T/\lVert u\rVert^2$, the reflection that fixes $u^\perp$; $\det H_u=-1$. A product of $j$
reflections fixes an intersection of $j$ hyperplanes, so it satisfies $\rk(R-I)\le j$. Every $R\in\Ort(n)$ is a product
of at most $n$ reflections (Cartan--Dieudonné). For the Euclidean form the bound is sharp (Scherk): if
$k=\rk(R-I)$, then $R$ is a product of exactly $k$ reflections and of no fewer. Hence $\det R=(-1)^{\rk(R-I)}$, and every
$R\in\SO(n)$ is a product of $\rk(R-I)$ reflections, an even number.
\end{auxresult}
\noindent\emph{Form used.} Scherk's theorem for a nondegenerate quadratic space adds two reflections when the moved space
$\mathrm{im}(R-I)$ is totally isotropic; a definite form has no nonzero isotropic vectors, so this case does not occur.
Companion: \bgref{D.11}.
\fbBusedin{\thmref{Theorem}{II.7}(a) ($\supseteq$, and the parity of $r$), \thmref{Theorem}{V.3}, corollary~5 (re-HRA).}

\begin{auxresult}{B.23}{the Cayley map}{\citep{xcayley1846quelques}; \citep{lezcanocasado2019cheap}}
Let $\Omega\in\so(n)$.
\begin{enumerate}[label=(\alph*)]
  \item $I-\Omega$ is invertible, since the eigenvalues of $\Omega$ are imaginary, and $C(\Omega)=(I+\Omega)(I-\Omega)^{-1}$
    lies in $\SO(n)$ and commutes with every polynomial in $\Omega$.
  \item $C(\Omega)-I=2\Omega(I-\Omega)^{-1}$, so $\rk(C(\Omega)-I)=\rk\Omega$; and $C(\Omega)=I+2\Omega+2\Omega^2+\dots$, so
    $dC_0=2\,\id$.
  \item $C$ is a diffeomorphism from $\so(n)$ onto the open dense subset $\{R\in\SO(n):-1\notin\mathrm{spec}\,R\}$, with
    inverse $R\mapsto(R-I)(R+I)^{-1}$. In particular it misses $\diag(-1,-1,1,\dots,1)$.
  \item Every element of $\SO(n)$ is a product of two elements of the image of $C$.
  \item (Neumann truncation.) $(I-Q)(I+Q+Q^2+Q^3)=I-Q^4$, so for skew $Q$
    \begin{equation*}
      R:=I+2Q+2Q^2+2Q^3+Q^4=C(Q)\,(I-Q^4),\qquad R\T R=(I-Q^4)^2 ,
    \end{equation*}
    because $Q^4$ is symmetric and commutes with the orthogonal matrix $C(Q)$; the singular values of $R$ are
    $\lvert1-\mu_k^4\rvert$, where $\pm i\mu_k$ are the eigenvalues of $Q$.
\end{enumerate}
\end{auxresult}
\noindent\emph{Reason for (d).} By \auxref{B.7}(b), $R=e^X$ with $X$ skew and eigenvalues $\pm i\varphi$,
$\varphi\in[0,\pi]$. Then $e^{X/2}$ has eigenvalues $e^{\pm i\varphi/2}\ne-1$, and $R=(e^{X/2})^2$. Companion:
\bgref{D.7}.
\fbBusedin{\thmref{Corollary}{II.6} (HF PEFT's default OFT), \thmref{Theorem}{II.7}(c) (the germ of the image at $W_0$)
and (d) (exact Cayley OFT never produces eigenvalue $-1$), \thmref{Theorem}{V.3}, corollaries~2 and~4, the LOFT example
of Exposé~II.B. The proof of \thmref{Proposition}{V.2} in Appendix~\ref{app:proofs} also cites it.}

\begin{auxresult}{B.24}{Lie subgroups, generation, and Yamabe's theorem}{\citep[Prop.~7.14, Thms~19.25,
19.26]{xlee2013introduction}; \citep{xyamabe1950arcwise}; \citep{xgoto1969arcwise}}
\begin{enumerate}[label=(\alph*)]
  \item In a connected Lie group every neighbourhood of the identity generates the group. A symmetric neighbourhood
    $U=U^{-1}$ therefore generates it as a monoid.
  \item For every Lie subalgebra $\mathfrak h$ of the Lie algebra of $G$ there is a unique connected Lie subgroup of $G$
    with Lie algebra $\mathfrak h$, namely the subgroup generated by $\exp\mathfrak h$. Lie subgroups are weakly
    embedded, so an inclusion $H\subseteq H'$ of Lie subgroups is smooth and gives $\mathrm{Lie}(H)\subseteq\mathrm{Lie}(H')$.
  \item (Yamabe.) Every arcwise connected subgroup of a Lie group is a Lie subgroup.
  \item Consequently, the subgroup generated by connected Lie subgroups $H_1,\dots,H_p$ of $G$ is the connected Lie
    subgroup $H$ with $\mathrm{Lie}(H)=\mathrm{Lie}\langle\mathfrak h_1,\dots,\mathfrak h_p\rangle$, the smallest Lie
    subalgebra containing every $\mathfrak h_i$.
\end{enumerate}
\end{auxresult}
\noindent\emph{Reason for (d).} The generated subgroup is arcwise connected, so it is a connected Lie subgroup $H$ by (c),
and $\mathrm{Lie}(H)\supseteq\mathfrak h_i$ for every $i$ by (b). The connected Lie subgroup $H'$ with Lie algebra
$\mathrm{Lie}\langle\mathfrak h_i\rangle$ contains every $H_i$, so $H\subseteq H'$, hence
$\mathrm{Lie}(H)=\mathrm{Lie}(H')$ and $H=H'$ by uniqueness. Companions: \bgref{D.4}, \bgref{D.8}, \bgref{D.9}.
\fbBusedin{\thmref{Theorem}{V.3}(b), (c) and corollary~4; the LOFT example of Exposé~II.B (supports revisited
across restarts).}

\begin{auxresult}{B.25}{first fundamental theorem for $\Ort(n)$; Witt's extension theorem}{\citep[Ch.~II]{xweyl1939classical};
\citep[Ch.~5]{xgoodman2009symmetry}; \citep[Ch.~III]{xartin1957geometric}}
\begin{enumerate}[label=(\alph*)]
  \item Every polynomial in the rows $w_1,\dots,w_m\in\R^n$ of $W$ that is invariant under $W\mapsto WR$, $R\in\Ort(n)$,
    is a polynomial in the inner products $\langle w_i,w_j\rangle$. Since the invariant polynomials of a compact group
    acting linearly separate its orbits, $WW\T$ is a complete invariant of the right $\Ort(n)$-action.
  \item (Witt, Euclidean case.) A linear isometry between two subspaces of $\R^n$ extends to an element of $\Ort(n)$.
\end{enumerate}
\end{auxresult}
\noindent\emph{Form used.} Part (b) gives the orbit statement of (a) directly: if $WW\T=W'W'^{\top}$, the map
$\sum_ic_iw_i\mapsto\sum_ic_iw_i'$ is a well-defined isometry of the row spaces, and its extension is $R\T$ with $W'=WR$.
Companion: \bgref{D.16}.
\fbBusedin{\thmref{Theorem}{II.5}(1), (2), \thmref{Corollary}{II.6}, \thmref{Corollary}{III.13} (left $\Ort(r)$-orbits of
full-rank $A_0$ are classified by $A_0\T A_0$).}

\begin{auxresult}{B.26}{the Kempf--Ness theorem in the real form used}{\citep{xkempf1979length};
\citep{xrichardson1990minimum}; \citep{xbirkes1971orbits}}
Let $V$ be a real inner-product space, $G\subseteq\GL(V)$ a real algebraic group that is closed under transpose,
$K=G\cap\Ort(V)$, and $\mathfrak p$ the symmetric elements of its Lie algebra. Define the moment map $\mu:V\to\mathfrak p$
by $\langle\mu(v),\xi\rangle=\tfrac12\tfrac{d}{dt}\big|_{t=0}\lVert e^{t\xi}v\rVert^2=\langle\xi v,v\rangle$. Then:
\begin{enumerate}[label=(\roman*)]
  \item $g=e$ is a critical point of $g\mapsto\lVert g\cdot v\rVert^2$ iff $\mu(v)=0$, and then $\lVert v\rVert\le
    \lVert g\cdot v\rVert$ for every $g\in G$;
  \item $G\cdot v$ is closed iff it meets $\mu^{-1}(0)$;
  \item if $G\cdot v$ is closed, $G\cdot v\cap\mu^{-1}(0)$ is a single $K$-orbit;
  \item the closure of every orbit contains exactly one closed orbit.
\end{enumerate}
\end{auxresult}
\noindent\emph{Form used.} For $\GL_r$ acting on $V=\R^{m\times r}\times\R^{r\times n}$ by $(Bg^{-1},gA)$ with the
Frobenius inner product, the image is closed under transpose (the adjoint of the action of $g$ is the action of
$g\T$), $K=\Ort(r)$, and $\tfrac12\tfrac{d}{dt}\big|_0(\lVert Be^{-t\xi}\rVert^2+\lVert e^{t\xi}A\rVert^2)=
\langle\xi,AA\T-B\T B\rangle$. So the moment map is $-\mu$ with $\mu=B\T B-AA\T$ as in \thmref{Proposition}{III.7}, its
zero set is the balanced slice, a full-rank orbit is closed (it is a whole fibre of $(B,A)\mapsto BA$), and by
\auxref{B.21} the least value of $\lVert B\rVert^2+\lVert A\rVert^2$ on it is $2\lVert BA\rVert_*$. Companions:
\bgref{D.29}, \bgref{D.30}.
\fbBusedin{\thmref{Proposition}{III.7}(a), \thmref{Corollary}{III.6} (decay as a Kempf--Ness potential),
\thmref{Proposition}{V.4}(a) (part (iv)).}

\begin{auxresult}{B.27}{closed orbits in the fibres of $(B,A)\mapsto BA$}{\citep[Ch.~5]{xgoodman2009symmetry};
\citep[Ch.~1]{xmumford1994geometric}; part (c) elementary}
\begin{enumerate}[label=(\alph*)]
  \item (First fundamental theorem for $\GL_r$.) The polynomial invariants of $\GL_r(\mathbb C)$ acting on
    $\mathbb C^{m\times r}\times\mathbb C^{r\times n}$ by $(Bg^{-1},gA)$ are generated by the entries of $BA$, the
    contractions of the rows of $B$ with the columns of $A$.
  \item (Reductive quotients.) For a complex reductive group acting on an affine variety, every fibre of the quotient
    map contains exactly one closed orbit, and it lies in the closure of every orbit in that fibre.
  \item (Real form used.) Fix $W\in\R^{m\times n}$ of rank $k\le r$. Among the real pairs with $BA=W$, those with
    $\rk B=\rk A=k$ form a single $\GL_r$-orbit $O_W$. It is closed and lies in the closure of every orbit in the fibre
    $\{BA=W\}$. Hence a continuous $\GL_r$-invariant function of $(B,A)$ is constant on each fibre of $(B,A)\mapsto BA$.
\end{enumerate}
\end{auxresult}
\noindent\emph{Reason for (c).} If $BA=W$ with $\rk B=\rk A=k$, then $\mathrm{col}\,B=\mathrm{col}\,W$ and
$\mathrm{row}\,A=\mathrm{row}\,W$; a gauge brings $B$ to $[U_0\ 0]$ for a fixed basis $U_0$ of $\mathrm{col}\,W$, which
fixes the first $k$ rows of $A$, and the stabiliser of $[U_0\ 0]$ clears the remaining rows, which lie in
$\mathrm{row}\,W$. The set $O_W=\{BA=W,\ \rk B\le k,\ \rk A\le k\}$ is closed. For a general pair in the fibre, split
$\R^r=C_1\oplus C_2\oplus C_3\oplus C_4$ with $\mathrm{im}\,A=C_1\oplus C_2$, $\ker B=C_2\oplus C_3$ and $\dim C_1=k$,
and let $g_t$ act by $1$, $t$, $1$, $t^{-1}$ on the four summands. As $t\to0$, $(Bg_t^{-1},g_tA)$ converges to a pair in
the fibre whose factors have rank at most $k$, so to a point of $O_W$. Companion: \bgref{D.31}.
\fbBusedin{\thmref{Proposition}{V.4}(a) (a continuous combination rule that is invariant under each adapter's $\GL_r$ is a
function of the adapters, also off the full-rank locus).}

% ================================================================================================
\subsubsection{Dynamics, optimisation and differentiation}

\begin{auxresult}{B.28}{existence and uniqueness under Carathéodory conditions}{\citep[Ch.~I, Thms~5.1 and
5.3]{xhale1980ordinary}}
Let $D\subseteq\R^{N+1}$ be open and let $F:D\to\R^N$ be measurable in $t$ for each fixed $x$ and continuous in $x$ for
each fixed $t$, with $\lvert F(t,x)\rvert\le m_U(t)$ on each compact $U\subseteq D$ for an integrable $m_U$. Then through
every $(t_0,x_0)\in D$ there is an absolutely continuous solution of $\dot x=F(t,x)$, holding for almost every $t$. If
moreover $\lvert F(t,x)-F(t,y)\rvert\le k_U(t)\lvert x-y\rvert$ on each compact $U$ for an integrable $k_U$, the solution
through each point is unique.
\end{auxresult}
\noindent\emph{Form used.} If $\gamma$ is measurable in $t$ and locally Lipschitz in $\theta$, with locally integrable
bounds and Lipschitz constants, and $\rho$ is of class $C^{1,1}_{\rm loc}$, then $F(t,q)=-\Lambda\,D\rho_q\T\gamma(t,\rho(q))$
satisfies these conditions, so the trajectory from $kq_0$ is unique.
\fbBusedin{\thmref{Theorem}{III.5}(ii) (the trajectory statement).}

\begin{auxresult}{B.29}{Noether's theorem for gradient flows}{\citep{kunin2020neural}; \citep[Prop.~5.1]{zhao2022symmetries};
\citep{marcotte2023abide}; \citep{du2018algorithmic}}
Let a Lie group act linearly on $Q=\R^N$ with $\rho(k\cdot q)=\rho(q)$, and let $\Lambda\succ0$ be constant. Along every
solution of $\dot q=-\Lambda\,D\rho_q\T\gamma(t)$, where $\gamma(t)$ is any cotangent signal, gradient or not, and for every
generator $\xi$ of the action with $\Lambda^{-1}\xi$ symmetric, the charge
\begin{equation*}
  J_\xi(q)=\tfrac12\langle q,\Lambda^{-1}\xi q\rangle
\end{equation*}
is constant. For LoRA with block rates $\eta_B,\eta_A$ and generators $\xi_X(B,A)=(BX,-XA)$, $X$ symmetric, the charges
are the entries of $\Phi=B\T B/\eta_B-AA\T/\eta_A$.
\end{auxresult}
\noindent\emph{Form used.} The cited sources treat gradient flows, \citet{zhao2022symmetries} with $\Lambda=I$; their
proof uses only $D\rho_q(\xi q)=0$, so it covers arbitrary signals $\gamma$ and constant $\Lambda$, which is how
\thmref{Theorem}{III.5}(i) states and proves it. Companions: \bgref{D.26}, \bgref{D.28}.
\fbBusedin{\thmref{Theorem}{III.5}(i), \thmref{Corollary}{III.6}, \thmref{Proposition}{III.7}(b),
\thmref{Theorem}{III.9} (the charge stays at $-\kappa I$).}

\begin{auxresult}{B.30}{Riemannian gradients on embedded submanifolds}{\citep[\S3.6]{xabsil2008optimization}}
Let $N\subseteq\R^D$ be an embedded submanifold with the metric induced by the Euclidean (Frobenius) inner product, and
$f$ smooth near $N$. The Riemannian gradient of $f|_N$ at $x$ is $P_{T_xN}\nabla f(x)$, the orthogonal projection of the
Euclidean gradient onto the tangent space, and the Riemannian gradient flow is $\dot x=-P_{T_xN}\nabla f(x)$.
\end{auxresult}
\noindent\emph{Form used.} For $N=W_0+\MM_r$ at $W=W_0+X$, with $U=\mathrm{col}\,X$ and $V=\mathrm{row}\,X$,
\auxref{B.5}(b) gives the projection $G\mapsto P_UG+GP_V-P_UGP_V$. Companion: \bgref{D.24}.
\fbBusedin{\thmref{Proposition}{III.4}(b) (LoRA-Pro's equivalent gradient).}

\begin{auxresult}{B.31}{Frobenius' theorem}{\citep[Prop.~19.3, Thms~19.12 and 19.21]{xlee2013introduction}}
Let $D$ be a smooth distribution of constant rank $k$ on an open set $U$.
\begin{enumerate}[label=(\alph*)]
  \item If $N$ is an integral manifold of $D$ ($T_\theta N=D_\theta$ for $\theta\in N$) and $X,Y$ are sections of $D$,
    then $[X,Y]_\theta\in D_\theta$ for every $\theta\in N$. Hence an integrable distribution is involutive.
  \item (Frobenius.) Every involutive distribution is completely integrable: each point has a chart in which $D$ is
    spanned by the first $k$ coordinate fields. In such a chart a curve with $\dot\theta(t)\in D_{\theta(t)}$ keeps its last
    coordinates constant, so it stays in one integral manifold.
  \item (Global Frobenius.) The maximal connected integral manifolds of an involutive distribution form a foliation of
    $U$.
\end{enumerate}
\end{auxresult}
\noindent Companion: \bgref{D.25}.
\fbBusedin{\thmref{Proposition}{IV.3}.}

\begin{auxresult}{B.32}{flatness and the Bures--Wasserstein metric}{\citep[Ch.~7]{xlee2018introduction};
\citep{xtakatsu2011wasserstein}; \citep{xbhatia2019bures}}
\begin{enumerate}[label=(\alph*)]
  \item Riemannian curvature is preserved by local isometries, and a Riemannian manifold is flat iff it is locally
    isometric to an open subset of a Euclidean space.
  \item On the positive definite $n\times n$ matrices, the Bures--Wasserstein metric is the quotient of the Euclidean
    (Frobenius) metric under $B\mapsto BB\T$, whose fibres are the orbits $B\mapsto BR$, $R\in\Ort(n)$; up to a constant
    factor fixed by normalisation, $B\mapsto BB\T$ is a Riemannian submersion onto it. Its sectional curvature is
    nonnegative, and for $n\ge2$ it is positive on some planes, so the metric is not flat.
\end{enumerate}
\end{auxresult}
\noindent\emph{Form used.} By (a), an immersion $\rho$ with $\dim Q=\rk D$ that is a local isometry from the flat
$(Q,\Lambda^{-1})$ onto a leaf forces the leaf metric to be flat. The map $\rho(B)=BB\T$ has positive-dimensional fibres,
so it is no immersion; its cometric $D\rho\,D\rho\T(G)=2(G\theta+\theta G)$ at $\theta=BB\T$ does not depend on the point
of the fibre, and by (b) it is the cometric of the curved Bures--Wasserstein metric.
\fbBusedin{\thmref{Proposition}{IV.3} (dynamic equivalence with a forward lens forces a flat leaf metric when
$\dim Q=\rk D$).}

\begin{auxresult}{B.33}{the cost of forward and reverse automatic differentiation}{\citep[Ch.~3 and
4]{xgriewank2008evaluating}}
Let $L:\R^N\to\R$ be evaluated by a straight-line program.
\begin{enumerate}[label=(\alph*)]
  \item Forward mode evaluates $L(\theta)$ together with one directional derivative, $DL(\theta)[p]$, which equals
    $\langle\nabla L(\theta),p\rangle$, at a cost bounded by a small constant multiple of the cost of evaluating $L$. It
    carries each tangent alongside its value and stores no trace.
  \item Reverse mode evaluates $\nabla L(\theta)$ at a cost bounded by a small constant multiple of the cost of evaluating
    $L$, independent of $N$, but needs the intermediate values of the evaluation, stored or recomputed, during the
    reverse sweep.
\end{enumerate}
\end{auxresult}
\noindent\emph{Form used.} The $k$ coordinates $\langle p_i,\nabla L(\theta)\rangle$ of $P\T\nabla L$ cost $k$ forward
passes and no reverse pass; $\nabla L$ itself by forward mode would cost $N$ of them.
\fbBusedin{\thmref{Observation}{IV.4}.}

\begin{auxresult}{B.34}{central differences and Gaussian smoothing}{\citep{xnesterov2017random}; part (a) by Taylor's
theorem}
Let $L:\R^N\to\R$ be twice differentiable with an $\ell_2$-Lipschitz Hessian (constant $L_2$), $\varepsilon>0$ and
$z\in\R^N$.
\begin{enumerate}[label=(\alph*)]
  \item $\Big\lvert\dfrac{L(\theta+\varepsilon z)-L(\theta-\varepsilon z)}{2\varepsilon}-\langle\nabla L(\theta),z\rangle
    \Big\rvert\le\dfrac{L_2}{6}\,\varepsilon^2\lVert z\rVert^3$.
  \item For $z\sim\mathcal N(0,I_N)$, $\mathbb E\big[\langle\nabla L(\theta),z\rangle z\big]=\nabla L(\theta)$, and for
    Lipschitz $L$, $\mathbb E\big[\tfrac{1}{2\varepsilon}(L(\theta+\varepsilon z)-L(\theta-\varepsilon z))\,z\big]=\nabla
    L_\varepsilon(\theta)$, the gradient of the Gaussian smoothing $L_\varepsilon(\theta)=\mathbb E_uL(\theta+\varepsilon
    u)$.
  \item Hence the central two-point estimator is unbiased for $\nabla L_\varepsilon$, and by (a) its bias relative to
    $\nabla L$ is at most $\tfrac{L_2}{6}\varepsilon^2\,\mathbb E\lVert z\rVert^4=\tfrac{L_2}{6}\varepsilon^2N(N+2)$.
\end{enumerate}
\end{auxresult}
\fbBusedin{\thmref{Observation}{IV.4} (MeZO's estimator), \thmref{Proposition}{IV.3} (instances: each MeZO step is a step
of a linear method up to its estimator). The proof of \thmref{Definition}{IV.1} in Appendix~\ref{app:proofs} also cites it.}

\begin{auxresult}{B.35}{extreme eigenvalues of Wishart-type matrices}{\citep{xmarchenko1967distribution};
\citep{xyin1988limit}; \citep{xbai1993limit}}
Let $A\in\R^{r\times n}$ have independent, identically distributed entries with mean $0$, variance $c/n$ and finite
fourth moment. As $r,n\to\infty$ with $r/n\to y\in(0,1]$, the empirical eigenvalue
distribution of $AA\T$ converges to the Marčenko--Pastur law supported on $[c(1-\sqrt y)^2,c(1+\sqrt y)^2]$, and the
largest and smallest eigenvalues converge almost surely to $c(1+\sqrt y)^2$ and $c(1-\sqrt y)^2$.
\end{auxresult}
\noindent\emph{Form used.} HF PEFT's Kaiming-uniform $A_0$ has entries $U(-n^{-1/2},n^{-1/2})$, so $c=1/3$, and the
eigenvalues of $A_0A_0\T$ fill a band of relative width $O(\sqrt{r/n})$ around $c$. No proof depends on this entry.
\fbBusedin{the commentary on \thmref{Theorem}{III.9} (exactness for Kaiming initialisations), Prediction~X.7.}

% ================================================================================================
\subsubsection{Elementary facts}

\begin{auxresult}{B.36}{three elementary facts}{elementary; part (c) in \citealp[Ch.~1]{xhorn2013matrix}}
\begin{enumerate}[label=(\alph*)]
  \item If $\sigma:\R\to\R$ is differentiable at $0$ and $\sigma(\ell t)=\ell\,\sigma(t)$ for all $\ell\ge0$ and
    $t\in\R$, then $\sigma$ is linear.
  \item A real vector space is not a finite union of proper subspaces. In particular, a hyperplane of $\R^d$ that is
    contained in a finite union of coordinate hyperplanes is one of them.
  \item For $X\in\R^{p\times q}$ and $Y\in\R^{q\times p}$, $XY$ and $YX$ have the same nonzero eigenvalues, with
    multiplicities.
\end{enumerate}
\end{auxresult}
\noindent\emph{Reasons.} (a) Taking $\ell=0$ gives $\sigma(0)=0$, and $\sigma(t)=t\,\sigma(1)$ for $t>0$ and
$\sigma(t)=-t\,\sigma(-1)$ for $t<0$; the one-sided derivatives at $0$ are $\sigma(1)$ and $-\sigma(-1)$, so
differentiability forces $\sigma(t)=t\,\sigma(1)$. ReLU and leaky ReLU satisfy the identity without being linear. (b)
If $V=V_1\cup\dots\cup V_k$ with $k$ minimal, pick $x\in V_1$ outside the other $V_i$ and $y\notin V_1$. No point $y+tx$
lies in $V_1$, so two of them, $t\ne t'$, lie in one $V_i$ with $i\ge2$, and then $x\in V_i$, a contradiction. The second
claim applies the first to $H=\bigcup_j(H\cap\{x_j=0\})$.
\fbBusedin{\thmref{Theorem}{VI.3} (rule R3: positively homogeneous activations), \thmref{Lemma}{VI.2}(b) (the kink
argument for ReLU), \thmref{Theorem}{III.9} (the nonzero eigenvalues of $T=\beta\beta\T$ are those of $\beta\T\beta$). The proof of \thmref{Corollary}{III.13} in Appendix~\ref{app:proofs} also cites it.}

% ================================================================================================
\subsubsection{Spectral projectors, Riemannian submanifolds and floating point}

\begin{auxresult}{B.37}{smooth spectral projectors}{\citep[Ch.~I, \S5 and Ch.~II, \S1]{xkato1995perturbation};
\citep[Ch.~III]{xbhatia1997matrix}}
Let $S:V\to\Sym(m)$ be smooth on an open set $V$ of a real vector space, with eigenvalues
$\lambda_1(W)\ge\dots\ge\lambda_m(W)$, and let $\lambda_r(W_*)>\lambda_{r+1}(W_*)$ at some $W_*\in V$. Then the gap
$\lambda_r>\lambda_{r+1}$ persists on a neighbourhood $V_*$ of $W_*$, and on $V_*$ the orthogonal projector $\Pi(W)$ onto the
sum of the eigenspaces of $\lambda_1(W),\dots,\lambda_r(W)$ is the Riesz projector
\[\Pi(W)=\frac1{2\pi i}\oint_\Gamma\big(z-S(W)\big)^{-1}\,dz\]
for a fixed contour $\Gamma$ in $\mathbb C$ that encloses $\lambda_1(W_*),\dots,\lambda_r(W_*)$ and no other eigenvalue of
$S(W_*)$. In particular $\Pi$ is smooth on $V_*$.
\end{auxresult}
\noindent\emph{Reason.} By Weyl's perturbation inequality the eigenvalues are $1$-Lipschitz in the operator norm, so
near $W_*$ the contour still separates the top $r$ eigenvalues and $z-S(W)$ is invertible on $\Gamma$. Matrix inversion
is smooth, so the integral is smooth in $W$, and at each $W$ it is the orthogonal projector onto the enclosed eigenspaces.
\fbBusedin{\thmref{Proposition}{IV.3} (step 9: the GaLore field $\Pi(W)$ is smooth where a singular-value gap exists).}

\begin{auxresult}{B.38}{isometries, totally geodesic submanifolds and the Gauss equation}{\citep[Ch.~4, 5 and
8]{xlee2018introduction}; \citep{xkobayashi1958fixed}}
Let $(M,g)$ be a Riemannian manifold.
\begin{enumerate}[label=(\alph*)]
  \item A geodesic is determined by its initial point and initial velocity, and an isometry of $(M,g)$ maps geodesics to
    geodesics.
  \item Let $F\subseteq M$ be an embedded submanifold with the induced metric. If every geodesic of $M$ with initial
    velocity tangent to $F$ stays in $F$ for small times ($F$ is \emph{totally geodesic}), the second fundamental form of
    $F$ vanishes, and by the Gauss equation the sectional curvature of $F$ on every $2$-plane $\Pi\subseteq T_xF$ equals
    the sectional curvature of $M$ on $\Pi$.
\end{enumerate}
\end{auxresult}
\noindent\emph{Form used.} Let $\tau$ be an isometry and $F$ an embedded submanifold, open in its fixed-point set, with
$d\tau=\id$ on $TF$. By (a) a geodesic $\gamma$ tangent to $F$ and $\tau\circ\gamma$ have the same initial data, so
$\tau\circ\gamma=\gamma$ and $\gamma$ stays in $F$ for small times. So $F$ is totally geodesic, and (b) applies.
\fbBusedin{\thmref{Proposition}{IV.3} (step 6: balanced full-rank LoRA induces a curved leaf metric).}

\begin{auxresult}{B.39}{exact scaling by powers of two in floating point}{\citep[Ch.~2]{xhigham2002accuracy}}
In binary floating-point arithmetic with correctly rounded operations (IEEE~754, rounding to nearest), let $c=2^k$ with
$k\in\mathbb Z$. If all operands, exact results and rounded results below stay in the normal range (no overflow and no
underflow), then $\mathrm{fl}(cx)=cx$, and for floating-point numbers $x,y$,
\[\mathrm{fl}(cx\pm cy)=c\,\mathrm{fl}(x\pm y),\quad\mathrm{fl}(cx\cdot y)=c\,\mathrm{fl}(xy),\quad
\mathrm{fl}(cx/y)=c\,\mathrm{fl}(x/y),\quad\mathrm{fl}\big(\sqrt{c^2x}\big)=c\,\mathrm{fl}\big(\sqrt x\big).\]
\end{auxresult}
\noindent\emph{Reason.} In the normal range, multiplication by $2^k$ only shifts the exponent, so it maps the set of
floating-point numbers onto itself and scales all distances between them by $c$. Rounding to the nearest floating-point
number, with its tie-breaking rule, therefore commutes with it, and each identity follows from
$\mathrm{fl}(cz)=c\,\mathrm{fl}(z)$ applied to the exact result $z$.
\fbBusedin{\thmref{Theorem}{III.10} (bit-exact trajectories when every block factor is a power of two).}
